\documentclass[11pt,letterpaper]{article}
\usepackage[margin=1in]{geometry}
\usepackage{graphicx} % Required for inserting images
\usepackage{tikz}
\usepackage{amsthm,amsmath,amssymb}
\usepackage{algorithm}
\usepackage{algpseudocode}
\usepackage{xcolor}
\usepackage{soul}
\usepackage{enumitem}

% Theorems
\newtheorem{theorem}{Theorem}[section]
\newtheorem{corollary}[theorem]{Corollary}
\newtheorem{lemma}[theorem]{Lemma}
\newtheorem{claim}[theorem]{Claim}
\theoremstyle{definition}
\newtheorem{definition}[theorem]{Definition}
\newtheorem{remark}[theorem]{Remark}
\newtheorem{observation}[theorem]{Observation}

\newcommand{\reals}{\mathbb{R}}
\newcommand{\Span}{\mathrm{span}}
\newcommand{\coreq}{\mathrm{coreq}}

\newcommand{\calA}{\mathcal{A}}
\newcommand{\calB}{\mathcal{B}}
\newcommand{\calC}{\mathcal{C}}
\newcommand{\calD}{\mathcal{D}}
\newcommand{\calE}{\mathcal{E}}
\newcommand{\calF}{\mathcal{F}}
\newcommand{\calG}{\mathcal{G}}
\newcommand{\calH}{\mathcal{H}}
\newcommand{\calI}{\mathcal{I}}
\newcommand{\calJ}{\mathcal{J}}
\newcommand{\calK}{\mathcal{K}}
\newcommand{\calL}{\mathcal{L}}
\newcommand{\calM}{\mathcal{M}}
\newcommand{\calN}{\mathcal{N}}
\newcommand{\calO}{\mathcal{O}}
\newcommand{\calP}{\mathcal{P}}
\newcommand{\calQ}{\mathcal{Q}}
\newcommand{\calR}{\mathcal{R}}
\newcommand{\calS}{\mathcal{S}}
\newcommand{\calT}{\mathcal{T}}
\newcommand{\calU}{\mathcal{U}}
\newcommand{\calV}{\mathcal{V}}
\newcommand{\calW}{\mathcal{W}}
\newcommand{\calX}{\mathcal{X}}
\newcommand{\calY}{\mathcal{Y}}
\newcommand{\calZ}{\mathcal{Z}}

\newcommand{\bfa}{\mathbf{a}}
\newcommand{\bfb}{\mathbf{b}}
\newcommand{\bfc}{\mathbf{c}}
\newcommand{\bfd}{\mathbf{d}}
\newcommand{\bfe}{\mathbf{e}}
\newcommand{\bff}{\mathbf{f}}
\newcommand{\bfg}{\mathbf{g}}
\newcommand{\bfh}{\mathbf{h}}
\newcommand{\bfi}{\mathbf{i}}
\newcommand{\bfj}{\mathbf{j}}
\newcommand{\bfk}{\mathbf{k}}
\newcommand{\bfl}{\mathbf{l}}
\newcommand{\bfm}{\mathbf{m}}
\newcommand{\bfn}{\mathbf{n}}
\newcommand{\bfo}{\mathbf{o}}
\newcommand{\bfp}{\mathbf{p}}
\newcommand{\bfq}{\mathbf{q}}
\newcommand{\bfr}{\mathbf{r}}
\newcommand{\bfs}{\mathbf{s}}
\newcommand{\bft}{\mathbf{t}}
\newcommand{\bfu}{\mathbf{u}}
\newcommand{\bfv}{\mathbf{v}}
\newcommand{\bfw}{\mathbf{w}}
\newcommand{\bfx}{\mathbf{x}}
\newcommand{\bfy}{\mathbf{y}}
\newcommand{\bfz}{\mathbf{z}}

\newcommand{\cost}{\mathrm{COST}}
\newcommand{\ALG}{\mathrm{ALG}}
\newcommand{\OPT}{\mathrm{OPT}}
\newcommand{\Flow}{\textsc{Flow}}
\newcommand{\ouralg}{\textsc{LineTracker}}

\newcommand{\EE}{\mathbb{E}}

\newcommand{\interval}[1]{\lfloor #1 \rceil}
\newcommand{\off}{\mathrm{off}}
\newcommand{\on}{\mathrm{on}}
\newcommand{\org}{\mathrm{org}}
\newcommand{\rel}{\mathrm{rel}}
\newcommand{\dis}[1]{#1}
\newcommand{\new}{\mathrm{new}}

\newcommand{\solution}{\medskip\noindent\textbf{Solution.}\par\medskip}

\newcommand{\Var}{\mathrm{Var}}

\usepackage[hidelinks]{hyperref}
\usepackage{cleveref}


\title{47-836: Network Optimization II \\ Homework 4}
\author{Poon Tze-Yang}
\date{November 2026}


\begin{document}

\maketitle

Make your solutions concise and formal. Your goal is to convince me that you know the solutions. You should typeset your solutions in Latex. Please submit on Canvas.

\begin{enumerate}[series=q]
\item Show that there is no ``additive'' $100$-approximation algorithm $A$ for TSP unless $\mathrm{P} = \mathrm{NP}$. We say that an algorithms is a $\alpha$ additive approximation if the solution produced by the algorithm has cost at most $\OPT(I) + \alpha$ for every instance $I$. (E.g.\ algorithm $A$ finds a Hamiltonian cycle with cost at most $\OPT(I) + 100$ for all instances $I$).
\end{enumerate}

\solution
We reduce Hamiltonian cycle to additively approximating TSP.
We are given an unweighted graph $G = (V,E)$ as an instance of Hamiltonian cycle.
Create the following weighted clique $G'$ on the same vertex set $V':=V$: if two vertices are connected in $G$, set the weight of the edge between them in $G'$ to be 0, else set it to 101.
If there exists a Hamiltonian cycle in $G$, then there exists a tour in $G'$ which does not use any edge of positive weight, and hence the optimal tour has cost 0.
If there does not exist a Hamiltonian cycle in $G$, then the optimal tour in $G'$ has cost at least 101.
If we could approximate TSP to within an additive factor of 100, running the above reduction on an instance of Hamiltonian cycle $G$, if we obtain a tour of cost at most 100, we know the instance is a yes instance, otherwise it is a no instance.
This would us decide Hamiltonian cycle in polynomial time (since the instance $G'$ can be constructed in poly time), which proves additive approximation of TSP to fixed factor is NP-hard.


\begin{enumerate}[resume=q]
\item Let $G$ be a bipartite graph with partitions $X$ and $Y$. Let $d(x)$ be the degree of a vertex $x$. If $d(x) \geq 1$ for all $x \in X$ and $d(x) \geq d(y)$ for all $xy \in E(G)$ where $x \in X$ and $y \in Y$ then show $G$ has a matching covering all vertices in $X$
\end{enumerate}

\solution
We verify the condition of Hall's theorem: for every subset $S\subseteq X$, the neighborhood of $S$, $N(S)$, has size at least $S$, i.e. $|S|\leq |N(S)|$.
Suppose for contradiction there exists $S\subseteq X$ such than $|S|>|N(S)|$.
Without loss of generality, select a minimal such $S$.
Then, every $B\subseteq N(S)$ has a neighborhood $N(B)$ such that $|N(B)|\geq |B|$ (If $|N(B)|<|B|$, then $|S\setminus N(B)|<|N(S)\setminus B|$, which contradicts minimality of $S$). 
Thus, there is a complete matching from $N(S)$ to $S$.
Let the image of this matching in $S$ be $C$.
Let $d(A)$ for a vertex set $A$ be the total degree of $A$.
By handshake lemma, $d(S) \leq d(N(S))$.
However, since there is a complete matching from $N(S)$ to $C$, the degree of each vertex in $N(S)$ is at most the degree of its partner in $C$, and $|C|<|S|$, we have that
\[
d(N(S)) \leq d(C) < d(S).
\]
This is a contradiction, which suffices to prove the claim.


\begin{enumerate}[resume=q]
\item Consider the following problem. You are given a bipartite graph $G$ with bipartition $X, Y$ and an integer $k$. Let $N(S)$ be the neighbors in $Y$ for a set of vertices $S \subseteq X$. You want to find a set $S \subseteq X$ such that $|S| \leq k$ and $|N(S)|$ is maximized. Find a constant approximation algorithm for this problem. Hint: consider a greedy algorithm that chooses vertices of $X$ one by one such that they cover the maximum number of uncovered neighbors in $Y$.
\end{enumerate}

\solution



\begin{enumerate}[resume=q]
\item Consider unweighted interval scheduling problem where we are given a collection of open intervals on the line and the goal is to find a maximum cardinality subset of nonoverlapping intervals. There is a simple greedy algorithm that gives an optimum solution. Consider the algorithm that orders the requests in increasing order of length and greedily selects them while maintaining feasibility. Show that this algorithm is a $1/2$-approximation.
\end{enumerate}

\solution



\begin{enumerate}[resume=q]
% sic: the sheet writes c(vq) in the triangle inequality; presumably c(vw) is meant.
\item Consider the Metric-TSP, we are given an undirected graph $G = (V, E)$ where for each $e \in E$ there is a cost $c(e) \geq 0$. The metric property ensures that $c(uw) \leq c(uv) + c(vq)$ for any $u, v, w \in V$ and that the graph is a complete graph. The goal is to find a Hamiltonian cycle of minimum length in $G$.

Consider the following algorithm which greedily builds a Hamiltonian Cycle. Let $u_i$ be the $i$th vertex the cycle visits and set $u_1$ to be any arbitrary vertex. In the $i$th iteration of the algorithm, the next vertex chosen to visit is the closest vertex to $u_i$ that is not in $\cup_{j \leq i} u_j$. That is, the closest unvisited vertex. The algorithm completes when all vertices are visited. Prove that this algorithm gives a $O(\log n)$-approximation for the Metric-TSP problem where $n = |V|$. (Hint: Break up the nodes based on their cost. Say nodes who cost between $2^i$ and $2^{i+1}$. Let his set be $X_i$. Can you lower bound this groups cost by the optimal solution?)
\end{enumerate}

\solution



\end{document}
